% ============================================================================
%  Viewpoint 1 — Corrected Version / 观点1 — 修正版
%  SCX Gauge Theory: Gauge Freedom in Mixture of Experts
%  Chinese + English bilingual
% ============================================================================
\documentclass[11pt,a4paper]{article}
\usepackage[UTF8]{ctex}
\usepackage{amsmath,amssymb,amsthm}
\usepackage{mathtools}
\usepackage{booktabs}
\usepackage[margin=2.5cm]{geometry}
\usepackage{xcolor}
\usepackage{hyperref}

\newtheorem{theorem}{Theorem}
\newtheorem{proposition}[theorem]{Proposition}
\newtheorem{definition}[theorem]{Definition}
\newtheorem{remark}[theorem]{Remark}

\newcommand{\R}{\mathbb{R}}
\newcommand{\gauge}{\mathbf{g}}

\title{\textbf{Viewpoint 1 — Corrected Statement}\\[0.3cm]
\large SCX Gauge Theory: What LayerNorm Does and Does Not Fix}
\author{}
\date{July 2, 2026 / 2026年7月2日}

\begin{document}
\maketitle

% ===================================================================
\section{Original Claim (Incorrect) / 原始观点 (有误)}
% ===================================================================

\begin{quote}
\textbf{Original Viewpoint 1:}\\[4pt]
Experts live in their own coordinate systems — this is not a bug but gauge freedom.\\
Each expert output carries an unobservable offset $g_m$.\\
Training loss does not sense it — the residual connection absorbs it.\\
Routers / auditors / comparators do not know about it.\\
They compare things that are incomparable.

\medskip
\textbf{原始观点1:}\\[4pt]
专家活在各自坐标系里——不是bug是规范自由度。\\
每个专家输出带不可观测偏移$g_m$。\\
训练损失不感知——残差连接吸收。\\
路由器/审计者/比较者不知道。\\
在比较不可比的东西。
\end{quote}

% ===================================================================
\section{Verdict / 判决}
% ===================================================================

\begin{center}
\fbox{\parbox{0.9\textwidth}{
\centering
\textbf{Verdict: LARGELY INCORRECT / 大部分错误}

The core intuition (parameter redundancy exists) has merit.\\
But every specific mechanism claim is wrong or requires heavy qualification.\\
LayerNorm eliminates constant offsets — \textbf{not} vector offsets.\\
Routers \textbf{can} detect gauge offsets.\\
Training loss \textbf{does} sense gauge offsets directly.
}}
\end{center}

% ===================================================================
\section{Corrected Viewpoint 1 / 修正后的观点1}
% ===================================================================

\begin{quote}
\fcolorbox{blue!30}{blue!5}{\parbox{0.95\textwidth}{
\textbf{Corrected Viewpoint 1:}\\[6pt]
Experts in a Mixture of Experts (MoE) operate in different coordinate systems:
each expert $m$ can carry an additive offset $\gauge_m \in \R^d$ in its output.
This offset is a \textbf{gauge degree of freedom} — a parameter redundancy rather
than a true observable.\\[4pt]
\textbf{However}, the claim that these offsets are invisible to the training
process, routers, and auditors is \textbf{incorrect}.\\[4pt]
\begin{enumerate}
    \item \textbf{LayerNorm:} In Pre-LN Transformers, LayerNorm eliminates the
    \textit{per-token constant} component of $\gauge_m$ (all dimensions shifted
    by the same scalar) via mean subtraction. It does \textbf{not} eliminate
    \textit{vector offsets} where different dimensions receive different shifts.
    The differential component $\gauge_m - \overline{\gauge_m}\cdot\mathbf{1}$
    \textbf{leaks through} LayerNorm and propagates to downstream layers.

    \item \textbf{Training loss:} The loss function directly senses $\gauge_m$.
    For MSE: $\mathcal{L}(y + \gauge_m, t) \neq \mathcal{L}(y, t)$.
    SGD compensates by adjusting learnable bias terms ($b \to b - \gauge_m$),
    which \textit{is} sensing — the loss landscape changes, and optimization
    responds.

    \item \textbf{Routers:} In Top-1 routing, $\gauge_m$ is \textbf{exactly
    recoverable} from observed outputs (error $<10^{-15}$). In multi-layer MoE,
    $\gauge_m$ from layer $l$ contaminates the input distribution of layer
    $l+1$, causing routing decisions to flip (empirically 5/8 tokens in a
    controlled experiment). Routers \textbf{do} know — their inputs are polluted.

    \item \textbf{Auditors:} Given sufficient samples where expert $m$ is
    activated, an auditor can recover $\gauge_m$ via ordinary least squares
    regression with error $<10^{-14}$. The offset is statistically identifiable.
\end{enumerate}

\medskip
\textbf{The real problem:} It is not that gauge offsets are invisible. It is
that existing architectures do not \textit{explicitly} model the gauge degrees
of freedom. Compensation happens \textbf{implicitly} — via SGD absorbing offsets
into bias terms, and LayerNorm partially suppressing them. This implicit
compensation degrades routing quality, pollutes inter-layer communication,
and makes expert outputs incomparable without proper gauge fixing.

\medskip
\noindent\rule{\textwidth}{0.5pt}

\medskip
\textbf{修正后的观点1:}\\[6pt]
专家在MoE中运行于不同的坐标系：每个专家$m$可以在其输出中携带加性偏移$\gauge_m \in \R^d$。
这个偏移是一个\textbf{规范自由度}——即参数冗余，而非真正的可观测量。\\[4pt]
\textbf{然而}，声称这些偏移对训练过程、路由器和审计者不可见是\textbf{错误的}。\\[4pt]
\begin{enumerate}
    \item \textbf{LayerNorm:} 在Pre-LN Transformer中，LayerNorm通过均值减法消除$\gauge_m$
    的\textit{逐token常数分量}（所有维度加相同标量），但\textbf{不}消除\textit{向量偏移}
    （不同维度加不同值）。差分分量$\gauge_m - \overline{\gauge_m}\cdot\mathbf{1}$
    \textbf{泄漏}通过LayerNorm传播到下游层。

    \item \textbf{训练损失:} 损失函数直接感知$\gauge_m$。
    对MSE: $\mathcal{L}(y + \gauge_m, t) \neq \mathcal{L}(y, t)$。
    SGD通过调整可学习偏置项来补偿（$b \to b - \gauge_m$），这\textit{就是}感知——
    损失景观发生了变化，优化过程做出了响应。

    \item \textbf{路由器:} 在Top-1路由中，$\gauge_m$可从观测输出中\textbf{精确恢复}
    （误差$<10^{-15}$）。在多层MoE中，第$l$层的$\gauge_m$污染第$l+1$层的输入分布，
    导致路由决策翻转（受控实验中5/8个token）。路由器\textbf{确实}知道——其输入被污染。

    \item \textbf{审计者:} 给定专家$m$被激活的足够样本，审计者可通过普通最小二乘回归
    恢复$\gauge_m$，误差$<10^{-14}$。该偏移是统计可辨识的。
\end{enumerate}

\medskip
\textbf{真正的问题是:} 并非规范偏移不可见，而是现有架构没有\textit{显式地}建模规范自由度。
补偿以\textit{隐式}方式发生——SGD将偏移吸收到偏置项中，LayerNorm部分抑制它们。
这种隐式补偿降低了路由质量，污染了层间通信，使专家输出在缺乏适当规范固定的情况下不可比较。
}}
\end{quote}

% ===================================================================
\section{Precise Analysis: What LayerNorm Does and Does Not Do}
% ===================================================================
\section{精确分析: LayerNorm 做了什么、没做什么}

\subsection{Mathematics of LayerNorm Gauge Absorption}
\subsection{LayerNorm 规范吸收的数学}

Consider a Pre-LN Transformer block where the output of layer $l$ passes
through LayerNorm before entering layer $l+1$:

\begin{equation}
    y^{(l)} = x^{(l)} + \text{FFN}^{(l)}\bigl(\text{LN}(x^{(l)})\bigr)
\end{equation}

If expert $m$ adds a gauge offset $\gauge_m$ to its output:

\begin{equation}
    y^{(l)} = x^{(l)} + \text{FFN}^{(l)}\bigl(\text{LN}(x^{(l)})\bigr) + \gauge_m
\end{equation}

The next layer's input is $\text{LN}(y^{(l)})$. For a vector $z \in \R^d$:

\begin{equation}
    \text{LN}(z)_i = \frac{z_i - \mu(z)}{\sigma(z)} \cdot \gamma_i + \beta_i,
    \quad \mu(z) = \frac{1}{d}\sum_{j=1}^d z_j
\end{equation}

Now consider two cases:

\medskip
\noindent\textbf{Case 1: Constant offset} $\gauge_m = c \cdot \mathbf{1}$ (all dimensions
shifted by the same scalar $c$).

\begin{align}
    \mu(y + c\mathbf{1}) &= \mu(y) + c \\
    y_i + c - \mu(y + c\mathbf{1}) &= y_i + c - (\mu(y) + c) = y_i - \mu(y) \\
    \sigma(y + c\mathbf{1}) &= \sigma(y)
\end{align}

Therefore $\text{LN}(y + c\mathbf{1}) = \text{LN}(y)$. The constant offset is
\textbf{completely eliminated}.

\medskip
\noindent\textbf{Case 2: Vector offset} $\gauge_m$ where $\gauge_{m,i}$ varies
across dimensions.

\begin{align}
    \mu(y + \gauge_m) &= \mu(y) + \overline{\gauge_m} \\
    (y_i + \gauge_{m,i}) - \mu(y + \gauge_m) &= (y_i - \mu(y)) + (\gauge_{m,i} - \overline{\gauge_m}) \\
    \sigma(y + \gauge_m) &\neq \sigma(y) \quad \text{(in general)}
\end{align}

\textbf{Only the mean component} $\overline{\gauge_m}$ is subtracted.
The \textbf{differential component} $\gauge_{m,i} - \overline{\gauge_m}$
survives LayerNorm and leaks to the next layer.

\medskip
\noindent\rule{\textwidth}{0.5pt}

\medskip
\noindent\textbf{情况1: 常数偏移} $\gauge_m = c \cdot \mathbf{1}$（所有维度加相同标量$c$）。

\begin{align}
    \mu(y + c\mathbf{1}) &= \mu(y) + c \\
    y_i + c - \mu(y + c\mathbf{1}) &= y_i + c - (\mu(y) + c) = y_i - \mu(y) \\
    \sigma(y + c\mathbf{1}) &= \sigma(y)
\end{align}

因此$\text{LN}(y + c\mathbf{1}) = \text{LN}(y)$。常数偏移被\textbf{完全消除}。

\medskip
\noindent\textbf{情况2: 向量偏移} $\gauge_m$，其中$\gauge_{m,i}$各维度值不同。

\begin{align}
    \mu(y + \gauge_m) &= \mu(y) + \overline{\gauge_m} \\
    (y_i + \gauge_{m,i}) - \mu(y + \gauge_m) &= (y_i - \mu(y)) + (\gauge_{m,i} - \overline{\gauge_m}) \\
    \sigma(y + \gauge_m) &\neq \sigma(y) \quad \text{(一般情况)}
\end{align}

\textbf{仅均值分量}$\overline{\gauge_m}$被消除。
\textbf{差分分量}$\gauge_{m,i} - \overline{\gauge_m}$穿透LayerNorm泄漏到下一层。

\subsection{Empirical Evidence / 实验证据}

\begin{table}[h]
\centering
\caption{Gauge offset suppression by LayerNorm}
\begin{tabular}{@{}lccc@{}}
\toprule
\textbf{Offset Type / 偏移类型} & \textbf{Input Norm} & \textbf{LN Residual} & \textbf{Leaks?} \\
\midrule
Constant $g = 3.0 \cdot \mathbf{1}$ & 12.00 & ${\sim}10^{-16}$ & No / 否 \\
Vector $g$ (norm 7.63)            & 7.63  & 4.26            & \textbf{Yes / 是} \\
\bottomrule
\end{tabular}
\end{table}

\begin{table}[h]
\centering
\caption{Summary of experimental evidence against original claims}
\begin{tabular}{@{}lll@{}}
\toprule
\textbf{Experiment / 实验} & \textbf{Result / 结果} & \textbf{Refutes / 反驳} \\
\midrule
Top-1 router $g_m$ recovery    & Exact (err $<10^{-15}$) & ``Router doesn't know'' \\
LN vector offset test          & Residual 4.26 leaks     & ``Residual absorbs'' \\
Direct loss sensing            & $\Delta\mathcal{L}=7.79$ & ``Loss doesn't sense'' \\
SGD bias absorption            & $b_{\text{no }g} \approx b_{\text{with }g}+g$ & Parameter redundancy $\checkmark$, not invisibility \\
Multi-layer route flip         & 5/8 tokens flip top-1   & ``Router doesn't know'' \\
Statistical $g_m$ separation   & Exact (err $<10^{-14}$) & ``Auditor doesn't know'' \\
\bottomrule
\end{tabular}
\end{table}

% ===================================================================
\section{Why Routers CAN Detect Gauge / 路由器为何能检测规范}
% ===================================================================

\subsection{Top-1 Routing: Direct Observation}
\subsection{Top-1 路由: 直接观测}

In Top-1 routing, exactly one expert is active per token. The MoE output is:

\begin{equation}
    y = \text{FFN}_{m^*}(x) + \gauge_{m^*}
\end{equation}

The gauge offset $\gauge_{m^*}$ appears \textbf{directly} in the output.
A downstream observer (or Layer $l+1$ router) sees $y$, which contains
$\gauge_{m^*}$ as an additive term. If the observer also knows $\text{FFN}_{m^*}(x)$
(or can approximate it), $\gauge_{m^*}$ is trivially recoverable by subtraction.

\medskip
\noindent\rule{\textwidth}{0.5pt}
在Top-1路由中，每个token恰好激活一个专家。MoE输出为:

\begin{equation}
    y = \text{FFN}_{m^*}(x) + \gauge_{m^*}
\end{equation}

规范偏移$\gauge_{m^*}$\textbf{直接}出现在输出中。下游观察者（或第$l+1$层路由器）
看到$y$，其中包含$\gauge_{m^*}$作为加性项。若观察者也知晓$\text{FFN}_{m^*}(x)$
（或可近似之），则$\gauge_{m^*}$可通过减法平凡地恢复。

\subsection{Multi-Layer Contamination / 多层污染}

In a multi-layer MoE, the Layer $l$ output feeds into Layer $l+1$'s router:

\begin{equation}
    r^{(l+1)} = \text{Router}^{(l+1)}\bigl(y^{(l)}\bigr)
              = \text{Router}^{(l+1)}\bigl(x^{(l)} + \text{FFN}^{(l)}(\cdot) + \gauge_m\bigr)
\end{equation}

The router computes $\alpha^{(l+1)}(x) = \text{softmax}(r^{(l+1)})$.
Since $\gauge_m$ shifts $r^{(l+1)}$, the softmax weights change:

\begin{equation}
    \alpha^{(l+1)}_k = \frac{\exp(r^{(l+1)}_k + \Delta_k)}{\sum_j \exp(r^{(l+1)}_j + \Delta_j)},
    \quad \Delta = W_{\text{router}} \cdot \gauge_m
\end{equation}

\textbf{Consequence:} $\gauge_m$ from layer $l$ changes the routing decisions
in layer $l+1$. This is direct evidence that routers \textbf{do} sense gauge offsets.

\medskip
\noindent\rule{\textwidth}{0.5pt}
在多层MoE中，第$l$层输出馈入第$l+1$层路由器:

\begin{equation}
    r^{(l+1)} = \text{Router}^{(l+1)}\bigl(y^{(l)}\bigr)
              = \text{Router}^{(l+1)}\bigl(x^{(l)} + \text{FFN}^{(l)}(\cdot) + \gauge_m\bigr)
\end{equation}

路由器计算$\alpha^{(l+1)}(x) = \text{softmax}(r^{(l+1)})$。
由于$\gauge_m$偏移了$r^{(l+1)}$，softmax权重发生变化:

\begin{equation}
    \alpha^{(l+1)}_k = \frac{\exp(r^{(l+1)}_k + \Delta_k)}{\sum_j \exp(r^{(l+1)}_j + \Delta_j)},
    \quad \Delta = W_{\text{router}} \cdot \gauge_m
\end{equation}

\textbf{结论:} 第$l$层的$\gauge_m$改变了第$l+1$层的路由决策。
这是路由器\textbf{确实}感知规范偏移的直接证据。

% ===================================================================
\section{How the Training Process Senses Gauge / 训练过程如何感知规范}
% ===================================================================

The claim ``training loss does not sense $\gauge_m$'' conflates two distinct
phenomena:

\begin{enumerate}
    \item \textbf{Instantaneous loss:} $\mathcal{L}(y + \gauge_m, t) \neq \mathcal{L}(y, t)$
    for any fixed $\gauge_m \neq 0$. The loss surface is \textbf{different}.
    Gradients $\partial\mathcal{L}/\partial\theta$ are shifted. Training
    \textbf{directly senses} $\gauge_m$.

    \item \textbf{Post-convergence equivalence:} If the expert has a learnable
    bias $b$, then the effective bias after training is $b + \gauge_m$.
    SGD can reach the same optimal loss value (up to optimization noise) by
    setting $b^*_{\text{with }g} = b^*_{\text{without }g} - \gauge_m$.
    This is \textbf{compensation}, not insensitivity.
\end{enumerate}

\noindent\rule{\textwidth}{0.5pt}

声称``训练损失不感知$\gauge_m$''混淆了两个不同的现象:

\begin{enumerate}
    \item \textbf{瞬时损失:} 对任意固定$\gauge_m \neq 0$，
    $\mathcal{L}(y + \gauge_m, t) \neq \mathcal{L}(y, t)$。
    损失景观是\textbf{不同的}。梯度$\partial\mathcal{L}/\partial\theta$被偏移。
    训练\textbf{直接感知}$\gauge_m$。

    \item \textbf{收敛后等效:} 若专家有可学习偏置$b$，则训练后的有效偏置为$b + \gauge_m$。
    SGD可以通过设定$b^*_{\text{带}g} = b^*_{\text{不带}g} - \gauge_m$
    达到相同的最优损失值（在优化噪声范围内）。这是\textbf{补偿}，而非不敏感。
\end{enumerate}

% ===================================================================
\section{Comparison with SCX Gauge Theory / 与SCX规范理论的对比}
% ===================================================================

\begin{table}[h]
\centering
\caption{SCX gauge theory vs.\ MoE gauge redundancy}
\begin{tabular}{@{}p{4cm}p{5.5cm}p{5.5cm}@{}}
\toprule
\textbf{Aspect / 方面} & \textbf{SCX Gauge Theory} & \textbf{MoE Gauge Redundancy} \\
\midrule
Gauge group / 规范群
    & $(\R, +)$ — single scalar shift
    & $(\R^d, +)^M$ — per-expert vector shift \\
Gauge fixing / 规范固定
    & \textbf{Active}: $\sum g_e = 0$ by convention
    & \textbf{Passive}: SGD + LayerNorm choose implicitly \\
Observability / 可观测性
    & True gauge invariance: physical observables unchanged
    & Pseudo-gauge: observables \textit{are} affected, compensated post-hoc \\
LayerNorm role / LN角色
    & Not applicable
    & Eliminates constant component only; vector component leaks \\
Router awareness / 路由器感知
    & Not applicable
    & Directly affected; routing decisions change \\
\bottomrule
\end{tabular}
\end{table}

The MoE situation is \textbf{not} a true gauge symmetry in the physical sense.
It is a \textbf{parameter redundancy} that the optimization process resolves
implicitly. Unlike SCX where $\sum g_e = 0$ is an active human convention,
MoE ``gauge fixing'' is a passive byproduct of gradient descent and
normalization layers.

\medskip
\noindent\rule{\textwidth}{0.5pt}
MoE的情况\textbf{不是}物理意义上的真正规范对称性。它是一种\textbf{参数冗余}，
优化过程隐式地解决它。与SCX中$\sum g_e = 0$是主动人为约定不同，
MoE的``规范固定''是梯度下降和归一化层的被动副产品。

% ===================================================================
\section{Required Qualifications for Any ``Gauge Freedom'' Claim}
% ===================================================================
\section{任何"规范自由度"声称所需的限定条件}

For the claim ``experts operate in different coordinate systems due to gauge
freedom'' to be accurately stated, the following qualifications are necessary:

\medskip
对于``专家因规范自由度运行于不同坐标系''的声称要准确表述，需要以下限定:

\begin{enumerate}
    \item \textbf{Learnable parameter condition / 可学习参数条件:}
    The expert must have learnable output biases capable of absorbing $\gauge_m$.
    Without learnable biases, $\gauge_m$ directly shifts the loss and is
    fully observable.

    \item \textbf{Offset type condition / 偏移类型条件:}
    In Pre-LN architectures, only \textit{per-token constant} offsets are
    eliminated by LayerNorm. \textit{Vector offsets} (different values per
    dimension) leak through. The claim ``residual connection absorbs'' is
    wrong — it is LayerNorm that does (partial) absorption.

    \item \textbf{Convergence condition / 收敛条件:}
    ``Loss doesn't sense $\gauge_m$'' is only true in the limit of perfect
    convergence where biases have fully absorbed the offset. During training,
    the loss surface differs and gradients are affected.

    \item \textbf{Single-layer condition / 单层条件:}
    ``Router doesn't know'' could only hold (weakly) for single-layer MoE
    where the router sees the original input $x$, not gauge-contaminated
    intermediate representations. In multi-layer MoE, Layer $l$ offsets
    propagate to Layer $l+1$ routing.

    \item \textbf{Information-limiting condition / 信息限制条件:}
    ``Auditor doesn't know'' requires the auditor to have insufficient samples
    to perform statistical separation. With adequate samples ($\gtrsim d$ per
    expert), $\gauge_m$ is recoverable via OLS regression with
    machine-precision accuracy.
\end{enumerate}

% ===================================================================
\section{Summary: The Correct Picture / 总结: 正确的图景}
% ===================================================================

\begin{center}
\fcolorbox{green!40}{green!5}{\parbox{0.92\textwidth}{
\textbf{The Correct Claim / 正确的声称:}\\[6pt]
Experts in MoE carry additive offsets $\gauge_m$ that constitute a form of
parameter redundancy. In Pre-LN Transformers, LayerNorm eliminates the
\textit{constant} (per-token mean) component of these offsets, but
\textit{vector} components (dimension-specific differences) leak through
the residual stream. Training senses these offsets directly and compensates
via bias terms — this is active sensing, not ignorance. Routers in multi-layer
MoE are directly affected: gauge offsets from one layer contaminate the input
distribution of the next, causing routing decisions to flip. Auditors with
sufficient samples can statistically recover the offsets with near-perfect
accuracy.\\[4pt]
\textbf{The real problem is not invisibility — it is that existing architectures
do not explicitly model the gauge degrees of freedom.} Implicit compensation
degrades routing quality and makes cross-expert comparisons unreliable.
Proper gauge fixing (explicitly constraining offsets, as in SCX's
$\sum g_e = 0$ convention) would make expert outputs comparable and improve
routing stability.

\medskip
\noindent\rule{\textwidth}{0.5pt}

\medskip
MoE中的专家携带加性偏移$\gauge_m$，构成一种参数冗余。在Pre-LN Transformer中，
LayerNorm消除这些偏移的\textit{常数}（逐token均值）分量，但\textit{向量}分量
（各维度不同的差值）通过残差流泄漏。训练直接感知这些偏移并通过偏置项进行补偿——
这是主动感知，而非无知。多层MoE中的路由器直接受影响：一层的规范偏移污染下一层的
输入分布，导致路由决策翻转。拥有足够样本的审计者可以统计性地恢复偏移，
精度接近机器精度。\\[4pt]
\textbf{真正的问题不是不可见性——而是现有架构没有显式地建模规范自由度。}
隐式补偿降低了路由质量，使跨专家比较不可靠。
适当的规范固定（显式约束偏移，如SCX的$\sum g_e = 0$约定）将使专家输出可比较，
并提高路由稳定性。
}}
\end{center}

% ===================================================================
\section*{Code Evidence / 代码证据}
% ===================================================================

All experimental claims are validated by \texttt{analysis\_viewpoint1.py}
(pure NumPy, no deep learning framework dependency). Key results:

\medskip
所有实验声称由\texttt{analysis\_viewpoint1.py}验证（纯NumPy，无深度学习框架依赖）。关键结果:

\begin{itemize}
    \item \textbf{PART 2:} Top-1 routing enables exact $g_m$ recovery (error $<10^{-15}$)
    \item \textbf{PART 3:} Constant offset eliminated by LN (residual ${\sim}10^{-16}$);
          vector offset leaks (residual 4.26)
    \item \textbf{PART 4:} Direct loss sensing: $\Delta\mathcal{L} = 7.79$
    \item \textbf{PART 5:} SGD absorbs $g$ into learnable bias (post-training difference norm 0.153)
    \item \textbf{PART 6:} Multi-layer contamination: 5/8 tokens flip top-1 routing
    \item \textbf{PART 7:} Statistical $g_m$ separation via OLS: error $<10^{-14}$
\end{itemize}

\end{document}
